\section{SeqGPLVM hyperparameters}

\begin{table}[htbp]
\centering
\footnotesize
\renewcommand{\arraystretch}{1.2}
\begin{threeparttable}
\caption{SeqGPLVM hyperparameters and initialization procedures in the final run}
\label{tab:seqgplvm_hparams_final}
\begin{tabularx}{\textwidth}{
    >{\raggedright\arraybackslash}p{0.24\textwidth}
    >{\raggedright\arraybackslash}p{0.28\textwidth}
    >{\raggedright\arraybackslash}X
}
\toprule
\textbf{Hyperparameter / process} & \textbf{Actual value in final run} & \textbf{Possible values / procedures available in code} \\
\midrule

\texttt{latent\_dim}
& 1
& Any integer. \\

\texttt{num\_inducing / n\_inducing\_x}
& 31
& Any integer. In the sweep, it is set as
\(\lfloor \sqrt{\texttt{train\_test\_split} \cdot n} \rfloor\). \\

\texttt{num\_inducing\_hidden}
& 13
& Any integer. \\

\texttt{treatment\_model}
& \texttt{BernoulliLikelihood}
& Any likelihood factory/class accepted by the model. The code explicitly handles \texttt{BernoulliLikelihood} and \texttt{GaussianLikelihood}; the Titsias branch requires \texttt{GaussianLikelihood}. \\

\texttt{z\_prior}
& \texttt{"normal"}
& \texttt{"normal"} or \texttt{"uniform"}. \\

\texttt{init\_z}
& \texttt{None}
& Either a user-supplied tensor/parameter or \texttt{None}. If \texttt{None}, the latent variables are initialized internally by random normal draws. \\

\texttt{z\_initializer}
& \texttt{"uniform"}
& \texttt{"uniform"} or \texttt{"normal"}. These map to different latent inducing-point initialization procedures. \\

\texttt{uniform\_halfwidth}
& 1
& Any positive float or \texttt{None}. Used by uniform latent inducing-point initialization. In the launched run, the sweep overwrites this value with \(a\). \\

\texttt{prior\_std}
& \texttt{None}
& Any positive float or \texttt{None}. Used by normal latent inducing-point initialization. \\

\texttt{learn\_inducing\_locations}
& \texttt{False}
& \texttt{True} or \texttt{False}. \\

\texttt{kernel} in \texttt{GPLVM}
& \texttt{"linear"}
& \texttt{"RBF"}, \texttt{"linear"}, \texttt{"rbf+linear"}, \texttt{"linear+rbf"}, \texttt{"rbf\_linear"}. \\

\texttt{optimize\_hyperparams["lr"]}
& \(10^{-2}\)
& Any numeric learning rate in the optimizer configuration. \\

\texttt{optimize\_hyperparams["num\_epochs"]}
& 6000
& Any positive integer. \\

\texttt{optimizer}
& \texttt{Adam}
& Only \texttt{torch.optim.Adam(model.parameters(), lr=...)} is implemented in the uploaded code. \\

\texttt{init\_z} initializer
& \(\texttt{torch.randn}(N,\texttt{latent\_dim})\)
& Two procedures are available: supply \texttt{init\_z} explicitly, or let the model initialize it by standard normal random draws. \\

Covariate inducing-point initialization
& Greedy farthest-point sampling via \texttt{farthest\_points(covs, n\_inducing\_x)}
& Active code path: greedy farthest-point sampling on observed non-missing covariates at each time point. The helper allows seeded or unseeded starts; the current call uses the default seeded behavior. \\

Latent Normal prior std
& 2.0
& Fixed at 2.0 in the current code when \texttt{z\_prior = "normal"}. Not exposed as a config hyperparameter. \\

GPLVM kernel priors
& For the effective \texttt{"linear"} kernel: \texttt{GammaPrior(2.0, 1.0)} on the linear variance, initialized at 1.0
& \texttt{"RBF"}: \texttt{GammaPrior(2.0,1.0)} on lengthscale and outputscale. \newline
\texttt{"linear"}: \texttt{GammaPrior(2.0,1.0)} on variance. \newline
\texttt{"rbf+linear"} variants: \texttt{GammaPrior(2.0,1.0)} on RBF lengthscale, RBF outputscale, and linear variance. \\

\bottomrule
\end{tabularx}

\begin{tablenotes}[flushleft]
\footnotesize
\item The ``final run'' refers to the actually launched run in the newest sweep script, i.e.\ the first selected configuration.
\item Although \texttt{uniform\_halfwidth} appears as 2.0 in the base training configuration, it is overwritten by the sweep with the current value of \(a\); hence the launched run uses \texttt{uniform\_halfwidth} \(=1\).
\item The effective \texttt{GPLVM} kernel is \texttt{"linear"} because the model constructor is called without explicitly overriding the default kernel.
\end{tablenotes}
\end{threeparttable}
\end{table}